\emph{Step 1 (the pair).} By Lemma \claimref{mass.de-sitter-scale-frozen-point}, $I_{\Lambda} = \nInf + \lambda\nO$ with $\lambda = \Lambda/6$ and $\ip{\nO}{\nInf} = 1$, so $Q(I_{\Lambda}) = 2\lambda = \Lambda/3$. The vector $W = \nInf - \lambda\nO$ satisfies $\ip{W}{I_{\Lambda}} = \lambda - \lambda = 0$ and $Q(W) = -2\lambda = -\Lambda/3$; in the frozen-point scale it is the tractor $(1, 0, -\Lambda/6)$, the value at the origin $x=0$ (for $\Lambda\neq0$) of the $W = V/Q(V)$ of Lemma \claimref{dark-energy.stelle-west-vector-frozen-point}, whose $e_{a}\wedge W$ are the model connection's translation bivectors.

\emph{Step 2 (stabiliser and square).} $e_{0}$ and $W$ are orthogonal, so $B_{\Lambda} = e_{0}W = e_{0}\wedge W$ is a $2$-blade. Both factors are orthogonal to $I_{\Lambda}$ (for $e_{0}$ because $I_{\Lambda}$ lies in the added plane), so $I_{\Lambda}$ anticommutes with each and commutes with their product. For orthogonal vectors $(e_{0}W)^{2} = -Q(e_{0})\,Q(W) = -(1)(-\Lambda/3) = \Lambda/3$.

\emph{Step 3 (the Killing field).} For a null point $X$ the vector $\xi = X\cdot B_{\Lambda} = \tfrac12(XB_{\Lambda} - B_{\Lambda}X) = \ip{X}{e_{0}}W - \ip{X}{W}e_{0}$ satisfies $\ip{X}{\xi} = 0$, so it is tangent to the cone. Because $B_{\Lambda}$ commutes with $I_{\Lambda}$, its flow fixes $I_{\Lambda}$ and so preserves the scale $\sigma = \ip{X}{I_{\Lambda}}$ that fixes the metric $g = \sigma^{-2}\eta$ below; a flow that is conformal for $\eta$ and preserves $\sigma$ is an isometry of $g$, so $\xi$ is a Killing field and not only a conformal Killing field. Writing $\xi = \delta X + aX$ with $\delta X$ the variation of the embedding along the induced chart velocity $\delta x$, and using $Q(X) = 0$, $\ip{X}{\delta X} = 0$ and $Q(\delta X) = Q(\delta x)$, one gets $Q(\xi) = Q(\delta x)$. With $\ip{X}{e_{0}} = x_{0}$ and $\ip{X}{W} = 1 + \lambda Q(x)/2$,
\[
Q(\xi) = x_{0}^{2}Q(W) + \bigl(1 + \tfrac{\lambda}{2}Q(x)\bigr)^{2} = \bigl(1 + \tfrac{\Lambda}{12}Q(x)\bigr)^{2} - \tfrac{\Lambda}{3}x_{0}^{2} .
\]
The metric is $g = \sigma^{-2}\eta$ with $\sigma = \ip{X}{I_{\Lambda}} = 1 - \lambda Q(x)/2$ (Lemma \claimref{mass.de-sitter-scale-frozen-point}). On $x = re_{1}$, $Q(x) = -r^{2}$, $\sigma = 1 + \lambda r^{2}/2$, the sphere of chart radius $r$ has areal radius $R = r/\sigma$, and
\[
g(\xi,\xi) = \frac{(1 - \lambda r^{2}/2)^{2}}{(1 + \lambda r^{2}/2)^{2}} = 1 - \frac{2\lambda r^{2}}{(1 + \lambda r^{2}/2)^{2}} = 1 - 2\lambda R^{2} = 1 - \frac{\Lambda}{3}R^{2},
\]
the static form $f(R) = 1 - \Lambda R^{2}/3$, with $g(\xi,\xi) = 1$ at the origin.

\emph{Step 4 (horizon, surface gravity, period).} For $\Lambda > 0$, $f(R_{h}) = 0$ at $R_{h} = \sqrt{3/\Lambda}$. By Fact \claimref{singularity.imported-surface-gravity-redshifted-acceleration} the surface gravity of a static Killing horizon is the limit of the redshifted acceleration of static observers, $|f'(R)|/2$, so $\kappa_{\Lambda} = \Lambda R_{h}/3 = \sqrt{\Lambda/3}$ and $\kappa_{\Lambda}^{2} = \Lambda/3 = B_{\Lambda}^{2}$. With $H = \sqrt{\Lambda/3}$ and $e_{W} = W/H$, $Q(e_{W}) = -1$ and $u = e_{0}e_{W}$ squares to $+1$: $B_{\Lambda} = Hu$ is $H$ times a unit boost in the plane of $e_{0}$ and $e_{W}$. On the slice $x_{0} = 0$, $x = re_{1}$, the velocity of Step 3 is exactly $\delta x = -(1 - \lambda r^{2}/2)\,e_{0}$ (there $\ip{X}{e_{0}} = 0$ and $\xi$ has no component in the added plane), so it vanishes at $x_{0} = 0$, $r^{2} = 2/\lambda$, that is $Q(x) = -12/\Lambda$, where $\sigma = 2$ and $R = r/2 = R_{h}$. Off the slice, matching the $\nO$, vector and $\nInf$ parts of $\xi = x_{0}W - (1 + \lambda Q(x)/2)e_{0}$ against $\delta X + aX$ with $\delta X = \delta x - \ip{x}{\delta x}\nInf$ gives $a = -\lambda x_{0}$ and
\[
\delta x = -\bigl(1 + \tfrac{\lambda}{2}Q(x)\bigr)e_{0} + \lambda x_{0}\, x .
\]
Writing $x = x_{0}e_{0} + y$ with $y$ spatial, for $\Lambda \neq 0$ this vanishes if and only if $\lambda x_{0}y = 0$ and $\lambda x_{0}^{2} = 1 + \lambda Q(x)/2$: either $x_{0} = 0$ and $Q(x) = -12/\Lambda$ (the bifurcation sphere above), or $y = 0$ and $x_{0}^{2} = 12/\Lambda$, the two points $x = \pm\sqrt{12/\Lambda}\,e_{0}$, where $\sigma = 1 - \lambda x_{0}^{2}/2 = 0$, outside the region the metric covers. Both need $\Lambda > 0$, and for $\Lambda = 0$, $\delta x = -e_{0}$ never vanishes. The velocity is quadratic in $x$, and at the bifurcation sphere its Jacobian is $\bigl(\begin{smallmatrix} 0 & H \\ H & 0\end{smallmatrix}\bigr)$ on the $(x_{0}, r)$ plane and zero on the sphere directions, eigenvalues $\pm H$, a coordinate-independent statement at a zero of a vector field. The argument of Proposition \claimref{singularity.euclidean-section-period} uses only $f(R_{h}) = 0$ and $f'(R_{h}) \neq 0$: in proper distance $\varrho$ from the horizon the $(\varrho, \tau)$ plane carries $d\varrho^{2} + \kappa_{\Lambda}^{2}\varrho^{2}d\tau^{2} + O(\varrho^{4})d\tau^{2}$, smooth at $\varrho = 0$ if and only if $\tau$ has period $2\pi/\kappa_{\Lambda}$.

\emph{Step 5 (the three signs).} Step 2 holds for every $\Lambda$. At $\Lambda = 0$, $W = \nInf$ and $B_{0} = e_{0}\nInf$, nonzero with square $-Q(e_{0})\,Q(\nInf) = 0$: a null rotation, the translation $e^{tN}X(x) = X(x - te_{0})$ of the flat model, with no fixed set and no horizon. For $\Lambda < 0$, $B_{\Lambda}^{2} = \Lambda/3 < 0$ and $\exp(tB_{\Lambda}/2)$ is periodic in $t$ with period $4\pi/\sqrt{-\Lambda/3}$, its conjugation action with period $2\pi/\sqrt{-\Lambda/3}$; $f(R) = 1 + |\Lambda|R^{2}/3$ has no zero.

\emph{Step 6 (reversal).} Under $Q \mapsto -Q$ the norms of $I_{\Lambda}$, $e_{0}$ and $W$ change sign, while $(e_{0}W)^{2} = -Q(e_{0})\,Q(W)$ is a product of two of them and does not.

The script \file{check.py} confirms every identity in exact arithmetic in $\Cl(2,4)$: (i) at eight values of $\Lambda$ of both signs and zero, and in the reversed form; (ii) on the product grid of three values of $\lambda$ and five values of each coordinate, which proves the polynomial identity since $Q(\xi)$ has degree at most two in $\lambda$ and four in each coordinate; the static-slice norm at six radii and four $\Lambda$; the horizon data at three rational $H$; and, at three bifurcation spheres, the vanishing of the velocity and its Jacobian, computed exactly by central differences since the velocity is quadratic.
